% Options for packages loaded elsewhere
\PassOptionsToPackage{unicode}{hyperref}
\PassOptionsToPackage{hyphens}{url}
\PassOptionsToPackage{dvipsnames,svgnames,x11names}{xcolor}
\documentclass[
  10pt,
]{article}
\usepackage{xcolor}
\usepackage[margin=1in]{geometry}
\usepackage{amsmath,amssymb}
\setcounter{secnumdepth}{-\maxdimen} % remove section numbering
\usepackage{iftex}
\ifPDFTeX
  \usepackage[T1]{fontenc}
  \usepackage[utf8]{inputenc}
  \usepackage{textcomp} % provide euro and other symbols
\else % if luatex or xetex
  \usepackage{unicode-math} % this also loads fontspec
  \defaultfontfeatures{Scale=MatchLowercase}
  \defaultfontfeatures[\rmfamily]{Ligatures=TeX,Scale=1}
\fi
\usepackage{lmodern}
\ifPDFTeX\else
  % xetex/luatex font selection
  \setmainfont[]{DejaVu Serif}
  \setsansfont[]{DejaVu Sans}
  \setmonofont[Scale=0.82]{DejaVu Sans Mono}
\fi
% Use upquote if available, for straight quotes in verbatim environments
\IfFileExists{upquote.sty}{\usepackage{upquote}}{}
\IfFileExists{microtype.sty}{% use microtype if available
  \usepackage[]{microtype}
  \UseMicrotypeSet[protrusion]{basicmath} % disable protrusion for tt fonts
}{}
\makeatletter
\@ifundefined{KOMAClassName}{% if non-KOMA class
  \IfFileExists{parskip.sty}{%
    \usepackage{parskip}
  }{% else
    \setlength{\parindent}{0pt}
    \setlength{\parskip}{6pt plus 2pt minus 1pt}}
}{% if KOMA class
  \KOMAoptions{parskip=half}}
\makeatother
\usepackage{longtable,booktabs,array}
\usepackage{caption}
\captionsetup[table]{skip=6pt}
\newcounter{none} % for unnumbered tables
\usepackage{calc} % for calculating minipage widths
% Correct order of tables after \paragraph or \subparagraph
\usepackage{etoolbox}
\makeatletter
\patchcmd\longtable{\par}{\if@noskipsec\mbox{}\fi\par}{}{}
\makeatother
% Allow footnotes in longtable head/foot
\IfFileExists{footnotehyper.sty}{\usepackage{footnotehyper}}{\usepackage{footnote}}
\makesavenoteenv{longtable}
\setlength{\emergencystretch}{3em} % prevent overfull lines
\providecommand{\tightlist}{%
  \setlength{\itemsep}{0pt}\setlength{\parskip}{0pt}}
\usepackage{bookmark}
\IfFileExists{xurl.sty}{\usepackage{xurl}}{} % add URL line breaks if available
\urlstyle{same}
% fallback for those not using the hyperref driver hyperxmp:
\makeatletter
\@ifundefined{xmpquote}{\newcommand{\xmpquote}[1]{#1}}{}
\makeatother
\hypersetup{
  pdftitle={Fault tolerance and recovery --- phixeron},
  colorlinks=true,
  linkcolor={blue},
  filecolor={Maroon},
  citecolor={Blue},
  urlcolor={blue},
  pdfcreator={LaTeX via pandoc}}

\title{Fault tolerance and recovery --- phixeron}
\author{}
\date{}

\begin{document}
\maketitle

{
\setcounter{tocdepth}{3}
\tableofcontents
}
\section{Fault tolerance and recovery ---
phixeron}\label{fault-tolerance-and-recovery-phixeron}

How this system survives node loss, leader failover, a stuck local
archive, a crashed gateway, and a lost network frame --- and how each
surviving/restarted part gets back to a correct state. This is a
description of what the code does today, not the aspirational superset
in \texttt{doc/0-overview.md} \ldots{}
\texttt{doc/6-detailed-architecture.md} --- cross-check against source
before trusting specifics there.

\subsection{0. The one governing invariant: no snapshots, full-log
replay}\label{the-one-governing-invariant-no-snapshots-full-log-replay}

Every recovery path in this document reduces to the same primitive:
\textbf{replay the sequenced log from \texttt{globalSeqNo} 1}.
\texttt{SequencerService.onTakeSnapshot} throws and \texttt{onStart}
refuses a snapshot image
(\texttt{src/main/java/org/limitless/phixeron/sequencer/SequencerService.java:295-300,544-550});
\texttt{clusterctl\ shutdown} uses Aeron's \texttt{ABORT} action, which
takes no snapshot, never \texttt{SHUTDOWN}
(\texttt{src/main/java/org/limitless/phixeron/tools/ClusterCtl.java:151-185}
--- \texttt{SHUTDOWN} snapshots first, which would silently break the
invariant below).

This is deliberate, not an oversight: \textbf{every node publishes and
records its own copy of the sequenced stream} (the ``tap'',
\texttt{aeron:ipc} stream 205 --- see §2), and a node restored from a
snapshot would hold a recording that starts wherever the snapshot did,
not at the beginning of the trading day. Full-log replay is what keeps
every node's tap recording a \emph{complete} copy of history, which is
what lets any node serve a cold-start or gap replay to its co-located
apps without depending on a peer. The cost is that recovery time and
archive size grow with uptime --- bounded in practice because the log is
scoped to one trading day (daily rollover), not unbounded
(\texttt{doc/todo.md}, \texttt{{[}{[}project-no-snapshots{]}{]}}
memory).

Everything downstream of the log --- \texttt{globalSeqNo}, the gateway
topology, FIX session sequence numbers, positions, reference data --- is
a \textbf{pure function of replaying that log}, so nothing needs its own
persistence or its own recovery procedure. That single property is why
the mechanisms below look structurally similar: kill something, let it
replay, done.

\subsection{1. Cluster / node-level fault tolerance (Java, Aeron Cluster
+
Raft)}\label{cluster-node-level-fault-tolerance-java-aeron-cluster-raft}

\subsubsection{1.1 Every node holds a complete, byte-identical copy of
history}\label{every-node-holds-a-complete-byte-identical-copy-of-history}

\texttt{SequencerService} runs on every cluster node --- leader and
follower alike --- and each one re-publishes every sequenced frame onto
its own node-local \texttt{aeron:ipc} tap
(\texttt{FEEDER\_CHANNEL}/\texttt{FEEDER\_STREAM\_ID},
\texttt{SequencerService.java:101-102}), which that node's co-located
Aeron Archive records. Because every node processes the same
Raft-committed log in the same order, the taps are byte-identical across
nodes --- there is no cross-node replication of the recording itself, no
leader-only archive, and no asymmetry between a follower's copy of
history and the leader's. The tap publication and its recording are
created once in \texttt{onStart} and live for the whole process,
continuous across leadership changes (\texttt{aeron:ipc} has no port to
collide on, unlike the retired UDP global stream this replaced), so a
node's recording is one continuous run spanning every leader tenure it
lived through.

Consequence: \textbf{any node a client is co-located with can serve full
history/gap replay}, and losing a node loses no history --- its peers
already hold an identical complete copy.

\subsubsection{1.2 Leader failover}\label{leader-failover}

Raft election is Aeron Cluster's own mechanism; phixeron's contribution
is what rides on top of it.
\texttt{SequencerService.onNewLeadershipTermEvent} fires on every node
on a new term and calls \texttt{applyLeadership}, which asks
\texttt{Sequencer.leadershipChanged} to synthesize a
\texttt{LeadershipChanged} frame --- de-duplicated against the leader
already on record, so a node re-observing its own term doesn't
double-emit (\texttt{Sequencer.java:392-407}). This frame is sequenced
and recorded exactly like any ingress message, so \textbf{every node's
tap recording --- not just the leader's --- carries a gap-free account
of every leadership change}, and a leader-only consumer (§3, §6) can
derive ``who is leader and since when'' purely from replaying the log.

On the C++ client side, \texttt{ClusterStreamSender} (the Aeron Cluster
ingress session state machine used by \texttt{FixGateway} and every
other cluster client) survives a leader failover \textbf{without losing
its cluster session}: a \texttt{NewLeaderEvent} swaps only the ingress
\texttt{Publication} to the new leader's endpoint, re-resolved out of
the event's member CSV --- the cluster session id and leadership term id
are updated in place, never re-created
(\texttt{src/main/cpp/org/limitless/phixeron/sequencer/ClusterStreamSender.hpp:632-689}).
\texttt{send()}'s retry loop pumps the egress control stream between
offer attempts specifically so an in-flight \texttt{NewLeaderEvent} can
land and swap the publication mid-spin --- a naive
\texttt{while(!offer)\ idle()} would deadlock, spinning on the dead
leader's publication while the poll that would revive it never runs
(\texttt{ClusterStreamSender.hpp:509-537}). A co-located client that
originally reached the leader over cheap IPC and failed over onto UDP
re-chases IPC if leadership later returns to its own member
(\texttt{ClusterStreamSender.hpp:642-672}).

A leader failover is explicitly \textbf{not} session loss --- the
fencing logic in §3 treats it as a transparent event, not a fault.

\subsubsection{1.3 A node that cannot record itself terminates
(self-fencing)}\label{a-node-that-cannot-record-itself-terminates-self-fencing}

\texttt{SequencerService.emit} is \emph{reliable}: it spins on the
tap-publication offer until it lands, because the recording is the
authoritative copy of history and a dropped frame would be an
unrecoverable gap (\texttt{SequencerService.java:640-688}). This can
only block on genuine local-archive back-pressure --- the tap's only
tethered subscriber is the recording itself, app replicas are untethered
--- but reliable is not the same as unbounded. \texttt{TapStallPolicy}
(pure, Aeron-free, unit-tested in isolation ---
\texttt{TapStallPolicy.java}) distinguishes an archive that is merely
slow (back-pressured but its recording position keeps advancing →
\texttt{CONTINUE}) from one that has stopped draining
(\texttt{FATAL\_NO\_PROGRESS} after 30s of zero progress) or gone away
entirely (\texttt{FATAL\_RECORDING\_GONE}). The same 1 Hz tick that
drives the cluster clock also runs \texttt{checkTapRecordingAlive}
(\texttt{SequencerService.java:480-486}), because a \emph{stopped}
recording doesn't back-pressure anything at all --- the tap's untethered
app subscribers keep it looking connected --- so liveness has to be
polled, not just inferred from back-pressure.

Either path calls \texttt{fatalTapFailure} → \texttt{fatalFailure},
which logs \texttt{FATAL:\ …\ terminating\ this\ node} and runs the
fatal handler wired by \texttt{SequencerServer}, exiting the process with
code \textbf{70} (\texttt{SequencerService.java:728-767}; documented
operator-facing in \texttt{doc/ops.md} ``A node that terminates
itself''). No cluster callback may signal failure by throwing instead:
\texttt{Image.boundedControlledPoll} has already advanced the log
position past the message before a thrown exception is caught, and
\texttt{AgentRunner} keeps the agent alive --- a throw here would
silently drop the frame and leave the node running with a hole in its
own recording, exactly the failure this whole mechanism exists to
prevent (class Javadoc, \texttt{SequencerService.java:66-74}). The same
reasoning bounds \texttt{scheduleTick}
(\texttt{SequencerService.java:517-538}): a consensus module that
refuses the cluster-clock timer for 30s continuous back-pressure is
wedged, not busy, and gets the same fatal treatment --- otherwise every
consumer's session clock silently stops advancing with no
operator-visible signal.

Consequence for the cluster: the remaining members hold identical,
complete recordings and keep quorum without the dead node (an election
moves leadership if it held it); \textbf{two nodes down is a quorum
loss}, not a repeat of the same event. Restarting the node is ordinary
full-log replay from \texttt{globalSeqNo} 1 (§0), rebuilding its tap
recording from scratch; if it exits 70 again immediately, the underlying
storage is still broken (start-up itself is bounded --- the recording
must go live within 5s, \texttt{TAP\_RECORDING\_START\_TIMEOUT\_NS},
\texttt{SequencerService.java:121}).

\subsubsection{1.4 Cluster shutdown /
restart}\label{cluster-shutdown-restart}

\texttt{clusterctl\ shutdown} is leader-gated: on the leader it
publishes an unsequenced \texttt{ClusterStopped} marker, best-effort
awaits its sequenced echo (so the log's last event before a planned stop
is always that marker), then calls Aeron's \texttt{ClusterTool.abort}
(\texttt{ClusterCtl.java:151-185}). \texttt{ABORT} --- not
\texttt{SHUTDOWN} --- is the only lifecycle action that terminates
without taking a snapshot, and \texttt{SequencerServer} wires
\texttt{ConsensusModule.Context.terminationHook} to a barrier so every
node still unwinds cleanly (closing its Archive and draining the tap
recording to disk) rather than being killed abruptly
(\texttt{doc/clusterctl.md}). \texttt{clusterctl\ start} is the
read-side complement: it publishes \texttt{ClusterStarted} and waits for
its sequenced echo, so an operator/script can confirm the cluster is
actually up (elected leader, ingress accepted) rather than merely that
processes launched.

\subsection{2. FIX gateway (C++) fault
tolerance}\label{fix-gateway-c-fault-tolerance}

\texttt{FixGateway} is a deliberately stateless proxy: authoritative FIX
session state (sequence numbers, session status) lives in the cluster
(§0), not in the gateway process, specifically so the gateway can crash
and restart without losing anything durable. Two mechanisms carry the
rest: fencing (stop serving before you're wrong) and standby promotion
(someone else takes over).

\subsubsection{\texorpdfstring{2.1 Fencing:
\texttt{closeSessions}}{2.1 Fencing: closeSessions}}\label{fencing-closesessions}

\texttt{FixGateway::closeSessions(reason)}
(\texttt{src/main/cpp/org/limitless/phixeron/fix/FixGateway.cpp:668-696})
is the single choke point that stops this instance serving TCP clients:
it snapshots every active FIX session's recoverable sequence state into
\texttt{m\_recoveredSessions}, releases every connectionId from the
CompID registry, hard-closes every client socket (\texttt{::close(fd)}
--- no graceful Logout is sent; the fence deliberately looks to the
cluster exactly like this process dying), and sets
\texttt{m\_gateOpen\ =\ false} --- which also closes the accept gate,
since \texttt{processSockets} only calls \texttt{acceptNewConnection()}
while \texttt{m\_gateOpen} is true. It is idempotent (no-ops if the gate
is already shut).

Three independent signals trigger it
(\texttt{FixGateway.cpp:296-325,\ 557-577}), and they are not equivalent
--- one of them is not a fault at all:

{\def\LTcaptype{none} % do not increment counter
\begin{longtable}[]{@{}
  >{\raggedright\arraybackslash}p{(\linewidth - 4\tabcolsep) * \real{0.3333}}
  >{\raggedright\arraybackslash}p{(\linewidth - 4\tabcolsep) * \real{0.3333}}
  >{\raggedright\arraybackslash}p{(\linewidth - 4\tabcolsep) * \real{0.3333}}@{}}
\toprule\noalign{}
\begin{minipage}[b]{\linewidth}\raggedright
Signal
\end{minipage} & \begin{minipage}[b]{\linewidth}\raggedright
Where
\end{minipage} & \begin{minipage}[b]{\linewidth}\raggedright
What happens after
\end{minipage} \\
\midrule\noalign{}
\endhead
\bottomrule\noalign{}
\endlastfoot
A \texttt{GatewayActive} names a \textbf{sibling} instance while this
one was active & \texttt{sequencedEvent},
\texttt{FixGateway.cpp:557-577} & Fences,
\texttt{m\_activated\ =\ false}, but \textbf{keeps running} --- drops to
standby, keeps following the tap, and can be re-promoted later. The
cluster session is untouched. \\
The cluster closes this gateway's ingress session
(\texttt{ClusterStreamSender::isSessionLost()}) & \texttt{doWork},
\texttt{FixGateway.cpp:315-319} & Fences and \textbf{exits the process}
(\texttt{running\ =\ false}) --- a lost session cannot be re-established
in-process (§2.3), so an instance in this state could never be promoted
again; fail closed by exiting rather than idling with no session. \\
No \texttt{Tick} from the co-located tap for
\texttt{TAP\_STALL\_TIMEOUT\_MS} (20s = 20× the 1 Hz tick period) &
\texttt{checkTapStall}, \texttt{FixGateway.cpp:352-369} & Closes its own
cluster session first, then fences and \textbf{exits}. Voluntary: rather
than sit ``connected'' behind a dead tap, it forces the same
session-loss path as the row above, which drives standby promotion on
the cluster side. \\
\end{longtable}
}

The tap-stall watchdog is gated on \texttt{m\_replayer.isCaughtUp()} so
an in-progress cold-start/gap replay never reads as a stall, and it
measures \textbf{monotonic wall-clock time}, not cluster-consensus time
--- deliberately, since consensus time is itself delivered by the very
\texttt{Tick} frames being watched for, so it would freeze along with a
stalled tap and never trip (\texttt{FixGateway.cpp:241-244}).

A \texttt{GatewayActive} naming \emph{this} instance while it was
standby is the mirror case: \texttt{m\_activated} flips true and the
accept gate can open once every other gate condition is met (§2.4) ---
no fence involved.

\subsubsection{2.2 Standby promotion}\label{standby-promotion}

Every logical gateway can have multiple instances (a \texttt{Gateway}
row per instance, ranked by \texttt{preferenceRank}, loaded via
BasicData --- §5). The \textbf{cluster}, not the gateway, decides which
instance is active, by naming a \texttt{gatewayId} in a sequenced
\texttt{GatewayActive} frame --- every instance (and every node's
Replayer-fed app) observes the same frame in the same order, so there is
never a window where two instances both believe they're active off
inconsistent information.

Two ways a \texttt{GatewayActive} is produced:

\begin{itemize}
\tightlist
\item
  \textbf{Automatic, on session close.} \texttt{Sequencer.sessionClosed}
  fires whenever a cluster session closes (crash, network loss, graceful
  shutdown --- Aeron Cluster reports all of these as
  \texttt{onSessionClose}) and checks whether the closed session was one
  an \emph{active} gateway had declared itself on via
  \texttt{GatewayStarted} (\texttt{Sequencer.java:174,435-445} ---
  deliberately keyed off \texttt{GatewayStarted}, not the routing
  \texttt{sourceId} on every message, because other clients legitimately
  echo a gateway's \texttt{sourceId} and an earlier version of this
  logic let an unrelated \texttt{OrderExecServer} restart promote a
  standby out from under a perfectly healthy primary). If so,
  \texttt{promotionTarget} picks the lowest-\texttt{preferenceRank}
  sibling sharing the same \texttt{gatewaySourceId} and emits
  \texttt{GatewayActive} naming it --- or emits nothing, fail-closed, if
  there is no sibling to hand over to rather than naming a nonexistent
  instance (\texttt{Sequencer.java:447-478}).
\item
  \textbf{Manual, via
  \texttt{clusterctl\ activate\ \textless{}gatewayId\textgreater{}}.}
  Publishes an unsequenced \texttt{GatewayActive} that flows through
  \texttt{Sequencer.sequenceMessage} like any other ingress message ---
  the same code path, no special-casing --- and waits for its sequenced
  echo, matched by \texttt{gatewayId}
  (\texttt{ClusterCtl.java:194-291}). This is the operator's lever for a
  planned failover.
\end{itemize}

A \textbf{bootstrap} activation also runs once per trading day: the
first \texttt{EndBasicData} (the completion marker of a reference-data
load, §5) triggers \texttt{Sequencer.pendingGatewayBootstrapActivation},
which names the rank-0 (\texttt{preferenceRank\ ==\ 0}) Gateway row ---
so exactly one instance opens its accept gate at cold start and every
sibling waits as a hot standby (\texttt{Sequencer.java:409-427}). A
bootstrap with no rank-0 row produces no frame --- fail closed rather
than guess.

\subsubsection{2.3 Recovering FIX session state after a
restart}\label{recovering-fix-session-state-after-a-restart}

Because \texttt{FixGateway} holds no durable state of its own, a
restarting instance rebuilds what it needs from two things: an in-memory
snapshot taken on the way \emph{out} (\texttt{closeSessions}, §2.1, and
the analogous \texttt{finalizeRecovery}), and lifecycle frames replayed
off the cluster stream on the way back \emph{in}.

\begin{itemize}
\tightlist
\item
  \textbf{On fencing/shutdown}, every active FIX session's
  \texttt{RecoveredSession} state --- next outgoing/expected sequence
  numbers plus the three-field inbound-gap state
  (\texttt{m\_inboundGapHighSeqNum},
  \texttt{m\_inboundGapRequestedThrough},
  \texttt{m\_inboundGapRunStart}, the fields \texttt{doc/gap.md}'s gap-1
  fix added specifically so a mid-gap restart doesn't resume with the
  gap silently reopened) --- is captured into
  \texttt{m\_recoveredSessions}, keyed by client CompID.
\item
  \textbf{On restart}, before the accept gate opens, the gateway replays
  \texttt{ClientConnected}/ \texttt{ClientDisconnected} lifecycle frames
  off the cluster stream to reconstruct placeholder ``recovering''
  connections (\texttt{FixConnection} built with \texttt{fd\ =\ -1}),
  then \texttt{finalizeRecovery} turns each surviving placeholder into a
  \texttt{RecoveredSession} entry (a crash never got to publish the
  matching \texttt{ClientDisconnected}, so the placeholder represents a
  session whose socket died with the old process) and bumps the
  next-connectionId counter past whatever was still pending recovery, so
  a freshly accepted TCP connection can never collide with one still
  being recovered (\texttt{FixGateway.cpp:634-664,698-734}).
\item
  \textbf{When a new TCP connection Logons} with a CompID matching an
  entry in \texttt{m\_recoveredSessions},
  \texttt{FixConnection::adoptRecoveredSession} restores the saved
  sequence/gap state onto the fresh connection and consumes (erases) the
  entry --- one-shot adoption, so the client resumes exactly where it
  left off rather than renegotiating from sequence 1.
\end{itemize}

\subsubsection{2.4 The accept gate and connect-once
semantics}\label{the-accept-gate-and-connect-once-semantics}

\texttt{processSockets} only opens the accept gate when \textbf{all} of
the following hold simultaneously: not already open,
\texttt{m\_activated} (named by a \texttt{GatewayActive}),
\texttt{m\_replayer.isCaughtUp()}, \texttt{m\_basicDataLoaded}, and
\texttt{m\_ingressSender.isConnected()}
(\texttt{FixGateway.cpp:384-391}) --- so a gateway structurally cannot
admit a FIX logon before its cluster session is live, its reference data
is loaded, and it has caught up on history.
\texttt{ClusterStreamSender::connect}/\texttt{connectColocated} runs
exactly once, from the constructor (\texttt{FixGateway.cpp:281-282}) ---
there is \textbf{no in-process reconnect loop}. Once a session is lost
(§2.1, row 2), the gateway exits rather than retrying; recovering
service is an external supervisor's job (or a sibling instance already
running as a promoted standby), not something this process attempts on
its own. This is a deliberate simplification: a session-loss retry loop
would have to re-derive whether it's still safe to be active, which the
promotion mechanism already decides externally and unambiguously.

\subsection{3. Stream recovery --- cold start, gaps, and the
live/history
split}\label{stream-recovery-cold-start-gaps-and-the-livehistory-split}

Every app that consumes the sequenced stream (the FIX gateway,
\texttt{OrderExecServer}, \texttt{BasicDataServer},
\texttt{fix\_test\_server}) uses the same split: read the co-located
\texttt{SequencerService}'s tap \textbf{directly and live} (untethered
\texttt{aeron:ipc?tether=false}, so a slow consumer is dropped rather
than back-pressuring the sequencer --- see §1.3 and
\texttt{doc/audit.md} S4), and ask the co-located
\texttt{ReplayerService} to fill in history on cold start or a detected
gap.

\subsubsection{\texorpdfstring{3.1 \texttt{ReplayerService} (Java, one
per
node)}{3.1 ReplayerService (Java, one per node)}}\label{replayerservice-java-one-per-node}

Off the live-delivery path entirely --- no app depends on it for
steady-state throughput, only for catching up. Before ever declaring
itself \texttt{ready}, it replays the first frame of its own oldest tap
recording and checks \texttt{globalSeqNo\ ==\ 1}; if that fails,
\texttt{ready} latches false for the process's lifetime
(\texttt{integrityFailed}) --- a deliberate refusal, because a first
frame that isn't 1 means this node's own recording is missing or
corrupted, and centralizing the check here means every app on the node
is told \texttt{ReplayUnavailable} instead of independently discovering
the same broken archive
(\texttt{src/main/java/org/limitless/phixeron/replayer/server/ReplayerService.java},
\texttt{checkReady}/ \texttt{peekFirstGlobalSeqNo}). An archive call
that throws mid-replay flips the service into a \texttt{stalled} state
--- retried at 1s intervals, answering requests \texttt{ReplayPending}
in the meantime --- without crashing the process or touching live
delivery, since live reads never go through this service.
\texttt{MAX\_CONCURRENT\_REPLAYS\ =\ 2} bounds archive-IO parallelism
(the only event that needs many concurrent replays is node
start/restart, when every co-located app cold-starts at once); a slot
freed by one client is handed to a waiting one immediately, with a 60s
idle-TTL as a backstop against a client that died mid-replay. Its own
control-plane replies are offered with a bounded spin and
\textbf{dropped} rather than blocked past that --- cheap, since the
requester just resends on a timer --- so one stuck app cannot couple
every other app's replay to it, the same untethered-drop philosophy as
the tap itself, applied to the control plane.

\subsubsection{\texorpdfstring{3.2 Client-side replay
(\texttt{ReplayerStreamReceiver},
C++)}{3.2 Client-side replay (ReplayerStreamReceiver, C++)}}\label{client-side-replay-replayerstreamreceiver-c}

Used by \texttt{FixGateway} and the other node-local app replicas. On
\texttt{start()}, it immediately requests a full walk from segment 0
(cold start). In steady state, a live-tap frame whose
\texttt{globalSeqNo} jumps ahead of the next expected value clears
\texttt{isCaughtUp()} and requests a \textbf{resume} at the last
dispatched position --- repairing just the hole rather than re-walking
the whole day's log. If the resumed replay's first frame doesn't match
the anchored \texttt{globalSeqNo} (meaning the active recording rotated
under the client --- the node it's reading from restarted), it falls
back to a full re-walk.

The one subtlety worth calling out for anyone touching this code:
\textbf{live-tap frames are retained, not dropped, while a walk/resume
is in flight.} A frame beyond the current hole is held in a bounded
pooled FIFO (\texttt{MAX\_MESSAGES\_FRAMES\ =\ 65536} /
\texttt{MAX\_MESSAGES\_BYTES\ =\ 16MB}, falling back to drop-and-rewalk
past that bound) and handed over the instant the in-flight replay
reaches the hole --- no residual gap at the seam. The commit history is
explicit about the cost of getting this wrong: dropping retained frames
(as this code did until 2026-08-05) meant every walk finished one
guaranteed hole short of live, and a single injected frame drop under a
300-message flood cost 6 full re-walks instead of 1 with retention.
\texttt{isCaughtUp()} is a state that can re-clear on a later gap, not a
one-time latch --- both \texttt{FixGateway}'s tap-stall watchdog (§2.1)
and leader-only emission gates (§6) depend on that, since an earlier
latch-forever bug meant a mid-recovery re-walk read as a stalled
sequencer and the gateway fenced itself out of a recovery it didn't
need. The very first frame a client ever dispatches must carry
\texttt{globalSeqNo\ ==\ 1} or the process aborts --- a hard invariant
that this node's recording reaches the start of the log.

The replay stream is subscribed \textbf{per episode and filtered to that
replay's own Aeron session id} (\texttt{openReplaySubscription}), never
held open between replays. That is a flow-control requirement, not
housekeeping. The \texttt{ReplayerService} answers every app on one
shared \texttt{aeron:ipc} stream (201), and an Aeron publication is
throttled by its slowest \emph{tethered} subscriber --- so a standing
subscription there (as this code had until 2026-08-07) made every idle
app a subscriber of every other app's replay, one that never polls,
since \texttt{poll()} only ever reads the image of its own session. Its
position sat at 0 forever, and the archive's replay wedged one
publication window in and never moved again: measured on a stalled cold
start, \texttt{pub-lmt} pinned at exactly 33,554,432 --- 32 MiB, half a
64 MB term --- with two peer \texttt{sub-pos} at 0 and the replay frozen
just past it at 34.6 MB, tens of MB short of its bound. Any cold start
needing more than \textasciitilde32 MiB of history therefore hung
\textbf{permanently}, which with no snapshots (§0) is what a restarted
replica faces after a few hundred thousand messages. Filtered by session
id, a replay publication has exactly one subscriber and no app can hold
back another's, while the stream stays tethered so a replay fragment is
still never silently dropped. This is the data-plane twin of the
control-plane coupling §3.1 avoids by dropping replies rather than
blocking on them; the live tap is the third case, and resolves it the
other way, by being untethered (§3 above, §1.3).

Nothing detected that hang either, which was the second half of the bug.
The resend timer only covers ``no \texttt{Replaying} yet''
(\texttt{m\_awaitingReplay}), and a bounded replay of a
\emph{still-recording} segment never closes its image on reaching its
bound --- completion there is detected by position, not by the image
closing --- so an image that is attached, open, and simply frozen had no
watchdog at all. A replay that makes no progress for
\texttt{REPLAY\_STALL\_TIMEOUT\_MS} (5s) now re-requests the same
segment verbatim: the same recovery already used when an image closes
\emph{short} of its bound (meaning the replay was stopped under the
client --- superseded, slot reclaimed, archive fault), extended to cover
a silent one, and a \texttt{Replaying} whose image never attaches at
all. 5s is far above any legitimate pause: the archive sustains tens of
MB/s into a local IPC replay, so a replay with anything left to serve is
never quiet for seconds, and a spurious fire only costs one segment
re-replayed.

\texttt{ClusterStreamClient} is the
archive-direct sibling used where there's no co-located Replayer
(\texttt{fix\_test\_server}, and \texttt{FixConnection}'s bounded
resend-recovery scan): it walks an archive's recorded segments for a
stream directly, replaying each historical segment fully and the last
(possibly still-recording) one open-ended, so the same image delivers
both historical and live messages without a subscription switch.

\subsection{\texorpdfstring{4. Exactly-once query replies across a
failover ---
\texttt{OutstandingQueries}}{4. Exactly-once query replies across a failover --- OutstandingQueries}}\label{exactly-once-query-replies-across-a-failover-outstandingqueries}

\texttt{PortfolioQueryRequest} handling (\texttt{OrderExecServer}) needs
a leader-only responder, but the leader can change mid-flight.
\texttt{OutstandingQueries.hpp} (templated, Aeron-free, unit-tested
standalone, mirroring the split between \texttt{Sequencer} and
\texttt{SequencerService}) keeps two separate notions of state:

\begin{itemize}
\tightlist
\item
  \textbf{Outstanding-ness is replicated} --- a request is outstanding
  from the moment it's sequenced until a matching \emph{sequenced} reply
  discharges it (\texttt{onRequest}/\texttt{onReply}), so every replica
  derives the same set of unanswered requests purely from the ordered
  log.
\item
  \textbf{Dispatch is local and advisory} --- ``a worker on this node is
  handling this one'' --- and gets cleared by \texttt{onNotLeader()} (a
  leadership change: the new leader must re-consider every
  still-outstanding request, including ones sequenced before it ever
  led) or \texttt{onReplyNotEmitted()} (a reply that was attempted but
  never reached the cluster, e.g.~a failed offer --- the request stays
  outstanding since the log is the source of truth).
\end{itemize}

A freshly promoted leader simply calls \texttt{dispatchUndispatched()}
against state every replica already holds identically --- no
special-cased failover recovery logic, no risk of answering a query
twice (discharge only ever happens via a sequenced reply) or losing one
(a reply that doesn't land leaves the request outstanding for the next
leader to pick up).

Dispatch runs in \textbf{insertion order}, which --- fed from the
sequenced stream --- is \texttt{globalSeqNo} order. That is load-bearing
rather than tidiness: these side effects are externally visible in the
order they are emitted (an \texttt{ExecutionReport} takes its outbound
FIX \texttt{MsgSeqNum} when the gateway frames it), so iterating the
underlying \texttt{unordered\_map} directly, as this did until
2026-08-07, handed a counterparty a burst of acks shuffled --- and
shuffled \emph{differently} per replica, since a rebuilt hash map's
iteration order is not a function of the log. Determinism across
replicas (§0) has to hold for emission order, not only for state.

Two more things ride the same machinery for the same reason.
\textbf{Venue acks}: every \texttt{NewOrderSingle} is outstanding from
the moment it is sequenced until its \texttt{ExecutionReport} is in the
log, keyed on the order's \texttt{globalSeqNo} --- which is also what
its \texttt{ExecID} names (\texttt{VenueExecId.hpp}). That replaced a
bare ``am I caught up and leader?'' test taken at the instant the order
was decoded, with no record kept, so an order arriving before this node
had processed its own promotion --- or delivered by a gap-repair replay,
during which \texttt{isCaughtUp()} is false throughout (§3.2) --- was
stepped past and acknowledged by nobody. Because the ack is a pure
function of the sequenced order (cluster timestamp included), a
re-emission after a failover is byte-identical to the one that may have
been lost, so the at-least-once seam yields exact duplicates that dedupe
on \texttt{ExecID} rather than two conflicting acks. \textbf{Rejected
queries} get a second, parallel
\texttt{OutstandingQueries\textless{}RejectedQuery\textgreater{}}
instance, so the canned ``queue full'' reply is retried and handed off
across a failover exactly like a real one instead of being
fire-and-forget.

\subsection{5. Reference data (BasicData)
recovery}\label{reference-data-basicdata-recovery}

\texttt{BasicDataServer}/\texttt{Gateways} treat reference data with the
same fault-tolerance shape as everything else: every node builds an
identical in-memory view by following its co-located tap, so losing a
node loses no reference data --- a neighbour already holds it
(\texttt{doc/basicdata-design.md} §0, mirroring \texttt{doc/audit.md}
S1). The producer role (reading the source-of-truth and publishing
\texttt{BasicData*} rows) is strictly leader-only and gated the same way
as query dispatch (§4) --- \texttt{isCaughtUp} plus ``is this node's
memberId the current leader'' --- so a failover doesn't produce two
competing loads; the new leader's replica simply opens its own upstream
connection and resumes.

Recovery deliberately has \textbf{no separate progress record}: a newly
promoted leader scans what's already in the log ---
\texttt{EndBasicData} seen → nothing to do; all sections complete but no
\texttt{EndBasicData} → emit it; otherwise resume the first incomplete
section from row 0 (\texttt{doc/basicdata-design.md} §4). One rule
covers every failure mode (DB error mid-section, plain crash, a failover
mid-load) uniformly, because it's derived from the log rather than
tracked separately from it --- the same principle as §0. Consumers
correspondingly commit a section only when its \texttt{remainingItems}
counter reaches 0, discarding any partial scratch buffer on a restart,
since this recovery rule can legitimately resend an in-flight section
after a failover.

\subsection{6. Operational tooling and
verification}\label{operational-tooling-and-verification}

\begin{itemize}
\tightlist
\item
  \textbf{\texttt{clusterctl}} (\texttt{doc/clusterctl.md}) ---
  \texttt{start}/\texttt{shutdown} bracket a run with sequenced markers
  so the log itself records ``the cluster was up between these two
  points''; \texttt{activate} is the manual standby-promotion lever
  (§2.2); \texttt{snapshot} is explicitly refused. See §1.4.
\item
  \textbf{Metrics} (\texttt{doc/ops.md}) ---
  \texttt{phixeron\_sequencer\_tap\_stalled} (latches at 1 when a node
  is about to terminate itself, §1.3),
  \texttt{phixeron\_sequencer\_gateway\_promotion\_total} (§2.2), the
  \texttt{phixeron\_replayer\_*} family (§3.1), and
  \texttt{phixeron\_node\_up} (an aggregator-synthesized per-node
  reachability gauge, independent of what else that node reports) give
  an operator the same signals this document describes, on a dashboard.
\item
  \textbf{\texttt{src/test/scripts/chaos-runner.sh}} --- randomized
  fault injection against a live 3-node cluster with a hot-standby
  gateway pair, replayable by seed. Injects:
  \texttt{fault\_kill\_leader}/\texttt{fault\_kill\_follower} (kill +
  restart a node, verifying quorum/leadership behave as above),
  \texttt{fault\_pause\_node} (\texttt{SIGSTOP} to simulate a GC pause),
  \texttt{fault\_tap\_drop} (one synthetic dropped live-tap frame,
  exercising §3.2's retained-message recovery), and
  \texttt{fault\_tap\_stall} (arms the same fault §1.3 self-terminates
  on, asserting the node actually exits rather than limping on with a
  dead recording). Between rounds it checks exactly one leader, a full
  FIX round trip against whichever gateway currently holds the accept
  gate, and that the live-tap consumer is still attached; at the end it
  freezes every member simultaneously and asserts every node's tap
  recording is gap-free, strictly monotone in \texttt{globalSeqNo}, and
  converged to the same high-water mark across all three nodes --- the
  strongest available proof that §1.1's ``byte-identical taps''
  invariant actually held for the run.
\item
  \textbf{\texttt{src/test/scripts/replay-bench.sh}} --- times a cold
  replica from launch to ``Caught up'' against a preloaded archive,
  optionally while load keeps arriving. It adds a fresh
  \texttt{OrderExecServer} beside a running cluster rather than
  restarting one (the launch script tears the cluster down when a child
  exits), so it walks the whole recording chain exactly as a restarted
  replica does. Written to chase §3.2's replay wedge and kept as the
  regression measurement for it: \texttt{replay-bench.sh\ 400000} builds
  \textasciitilde70 MB of history --- the case that used to hang forever
  and now converges in well under a second --- and a run reporting
  \texttt{NEVER\ CAUGHT\ UP} is that class of bug rather than a slow
  machine. This is also the measurement that matters for
  \texttt{fault\_kill\_leader} above, since a restarted replica has to
  catch up inside the harness's probe window; the replay wedge is what
  made those rounds fail.
\end{itemize}

\subsection{7. What this does not cover}\label{what-this-does-not-cover}

\begin{itemize}
\tightlist
\item
  \textbf{No pre-trade risk gating and no matching engine} --- a
  failover-safe query responder (§4) is not the same as a durable
  \emph{order acceptance} decision; see \texttt{doc/todo.md} items 1--2.
\item
  \textbf{No edge authentication} --- CompID validation only; see
  \texttt{doc/todo.md} item 3. Orthogonal to recovery, but relevant if a
  ``fault'' is ever adversarial rather than accidental.
\item
  \textbf{\texttt{aeronmd} itself and network partition/latency faults}
  are not exercised by \texttt{chaos-runner.sh} --- noted there as
  unwired seams (killing the media driver is destructive to co-located
  C++ clients; \texttt{tc\ netem}/\texttt{dnctl} isn't wired up on the
  macOS dev host).
\item
  \textbf{Single-node dev launches} have no failover to exercise at all
  --- the mechanisms above only engage with 2+ cluster members.
\end{itemize}

\end{document}
